\documentclass[openany]{book}
\usepackage[utf8]{inputenc}
\usepackage[T1]{fontenc}
\usepackage{blindtext}
\usepackage{graphicx}
\usepackage{hyperref}
\title{Testing tex4ebook to generate ePub: The Bare Basics}
\author{LianTze Lim}

\begin{document}

\maketitle

\tableofcontents

This is very experimental, and all caveats (and more) mentioned at the \href{https://www.overleaf.com/latex/examples/testing-html-export-with-make4ht/rqcknjvwyyry}{`make4ht hackety hack' example} applies to this example as well.

When compilation is complete, use the steps at \href{https://www.overleaf.com/learn/how-to/View_generated_files}{\emph{View generated files}} to download the generated .epub file.

\begin{figure}[hbt!]
\centering
\includegraphics[width=3cm]{example-image.png}
\includegraphics[width=3cm]{overleaf.pdf}
\caption{Hello!}
\label{fig:hello}
\end{figure}

This is a cross-reference to \ref{fig:hello}. Here are test citations: \cite{NTLKProject2015,Bond2014}

\begin{table}[hbt!]
\centering
\begin{tabular}{l|r}
Item & Quantity \\\hline
Widgets & 42 \\
Gadgets & 13
\end{tabular}
\caption{An example table.}\label{tab:widgets}
\end{table}

\blindmathpaper

\Blinddocument

\bibliographystyle{IEEEtran}
\bibliography{references}

\end{document}
